% Compile with XeLaTeX or LuaLaTeX for Avenir/Avenir Next.
% In VS Code LaTeX Workshop, choose recipe: latexmk (xelatex) or lualatex.

\documentclass[aspectratio=169,11pt]{beamer}

% ------------------------------------------------------------
% Theme foundation
% ------------------------------------------------------------
\usetheme{Antibes}
\useoutertheme{miniframes}      % top progress circles
\useinnertheme{rectangles}
\setbeamertemplate{navigation symbols}{} % no navigation icons near footer

% ------------------------------------------------------------
% Font setup: Avenir if available, graceful fallback otherwise
% ------------------------------------------------------------
\usepackage{iftex}
\ifPDFTeX
  % Fallback for accidental pdfLaTeX compilation.
  % For the intended look, compile with XeLaTeX or LuaLaTeX.
  \usepackage[T1]{fontenc}
  \usepackage[utf8]{inputenc}
  \usepackage{helvet}
  \renewcommand{\familydefault}{\sfdefault}
\else
  \usepackage{fontspec}
  \IfFontExistsTF{Avenir Next}{%
    \setsansfont{Avenir Next}[
      UprightFont = *,
      BoldFont = * Demi Bold,
      ItalicFont = * Italic,
      BoldItalicFont = * Demi Bold Italic
    ]
  }{%
    \IfFontExistsTF{Avenir}{%
      \setsansfont{Avenir}
    }{%
      % TeX Live fallback with a similar clean sans feel.
      \setsansfont{TeX Gyre Heros}
    }%
  }%
  \renewcommand{\familydefault}{\sfdefault}
\fi

% Keep math in a traditional serif math font for readability.
\usefonttheme{professionalfonts}
\usefonttheme[onlymath]{serif}


% ------------------------------------------------------------
% Packages
% ------------------------------------------------------------
\usepackage{amsmath,amssymb,mathtools}
\usepackage{graphicx}
\usepackage{xcolor}
\usepackage{tikz}
\usepackage{hyperref}
\usetikzlibrary{calc}

% ------------------------------------------------------------
% Color palette
% ------------------------------------------------------------
\definecolor{TwilightIvory}{HTML}{F6F1EA}
\definecolor{TwilightIvorySoft}{HTML}{FBF8F3}
\definecolor{TwilightCarbon}{HTML}{111111}
\definecolor{TwilightBlack}{HTML}{080808}
\definecolor{TwilightInk}{HTML}{050506}

\definecolor{TwilightDarkViolet}{HTML}{1E1533}
\definecolor{TwilightPurple}{HTML}{43305C}
\definecolor{TwilightPlum}{HTML}{2A1A42}
\definecolor{TwilightMulberry}{HTML}{7A4A63}

\definecolor{TwilightDeepTeal}{HTML}{0D7C7B}
\definecolor{TwilightAqua}{HTML}{A8E5DA}
\definecolor{TwilightBlush}{HTML}{F0D2DF}
\definecolor{TwilightRose}{HTML}{D8A4BF}

% ------------------------------------------------------------
% Global colors
% ------------------------------------------------------------
\setbeamercolor{normal text}{fg=TwilightCarbon,bg=TwilightIvory}
\setbeamercolor{background canvas}{bg=TwilightIvory}
\setbeamercolor{structure}{fg=TwilightDeepTeal}
\setbeamercolor{alerted text}{fg=TwilightMulberry}
\setbeamercolor{example text}{fg=TwilightDeepTeal}

\setbeamercolor{title}{fg=TwilightDarkViolet,bg=TwilightIvory}
\setbeamercolor{subtitle}{fg=TwilightPurple,bg=TwilightIvory}
\setbeamercolor{author}{fg=TwilightDeepTeal,bg=TwilightIvory}
\setbeamercolor{institute}{fg=TwilightCarbon,bg=TwilightIvory}
\setbeamercolor{date}{fg=TwilightCarbon,bg=TwilightIvory}

\setbeamercolor{frametitle}{fg=TwilightIvory,bg=TwilightDarkViolet}
\setbeamercolor{framesubtitle}{fg=TwilightAqua,bg=TwilightDarkViolet}

% ------------------------------------------------------------
% Title Slide Image
% ------------------------------------------------------------

% ------------------------------------------------------------
% Header bars: black, dark violet, purple descending from top
% ------------------------------------------------------------
\setbeamercolor{headline top}{fg=TwilightIvory,bg=TwilightBlack}
\setbeamercolor{headline middle}{fg=TwilightAqua,bg=TwilightDarkViolet}
\setbeamercolor{headline bottom}{fg=TwilightBlush,bg=TwilightPurple}

\setbeamercolor{section in head/foot}{fg=TwilightIvory,bg=TwilightBlack}
\setbeamercolor{subsection in head/foot}{fg=TwilightAqua,bg=TwilightDarkViolet}
\setbeamercolor{section in head/foot shaded}{fg=TwilightIvory!45!TwilightBlack,bg=TwilightBlack}
\setbeamercolor{subsection in head/foot shaded}{fg=TwilightBlush,bg=TwilightDarkViolet}



% Custom three-bar headline.  The bottom bar carries the miniframe progress meter.
\setbeamertemplate{headline}{%
  \leavevmode%
  \begin{beamercolorbox}[wd=\paperwidth,ht=2.25ex,dp=.75ex,leftskip=1.1em,rightskip=1.1em]{headline top}%
    \usebeamerfont{section in head/foot}
    \progressbar

  \end{beamercolorbox}%
  \begin{beamercolorbox}[wd=\paperwidth,ht=2.15ex,dp=.65ex,leftskip=1.1em,rightskip=1.1em]{headline middle}%
    \usebeamerfont{subsection in head/foot}\insertsubsectionnavigationhorizontal{\paperwidth}{}{}%
  \end{beamercolorbox}%
  \begin{beamercolorbox}[wd=\paperwidth,ht=2.05ex,dp=.6ex,leftskip=1.1em,rightskip=1.1em]{headline bottom}%
    \insertnavigation{\paperwidth}%
  \end{beamercolorbox}%
}

% ------------------------------------------------------------
% Footer: light cream/aqua/pink text, no navigation symbols
% ------------------------------------------------------------
\setbeamercolor{footline title}{fg=TwilightIvory,bg=TwilightBlack}
\setbeamercolor{footline author}{fg=TwilightAqua,bg=TwilightBlack}
\setbeamercolor{footline page}{fg=TwilightIvory,bg=TwilightBlack}

\setbeamertemplate{footline}{%
  \leavevmode%
  \hbox{%
    \begin{beamercolorbox}[wd=.46\paperwidth,ht=2.7ex,dp=1ex,leftskip=1em]{footline title}%
      \usebeamerfont{footline}\insertshorttitle
    \end{beamercolorbox}%
    \begin{beamercolorbox}[wd=.34\paperwidth,ht=2.7ex,dp=1ex,leftskip=.5em]{footline author}%
      \usebeamerfont{footline}\insertshortauthor
    \end{beamercolorbox}%
    \begin{beamercolorbox}[wd=.20\paperwidth,ht=2.7ex,dp=1ex,rightskip=1em plus1fil]{footline page}%
      \usebeamerfont{footline}\hfill\insertframenumber{} / \inserttotalframenumber
    \end{beamercolorbox}%
  }%
}

% ------------------------------------------------------------
% Blocks
% ------------------------------------------------------------
\setbeamercolor{block title}{fg=TwilightIvory,bg=TwilightDarkViolet}
\setbeamercolor{block body}{fg=TwilightCarbon,bg=TwilightIvorySoft}

\setbeamercolor{block title alerted}{fg=TwilightIvory,bg=TwilightMulberry}
\setbeamercolor{block body alerted}{fg=TwilightCarbon,bg=TwilightBlush!28!TwilightIvorySoft}

\setbeamercolor{block title example}{fg=TwilightIvory,bg=TwilightDeepTeal}
\setbeamercolor{block body example}{fg=TwilightCarbon,bg=TwilightAqua!20!TwilightIvorySoft}

% ------------------------------------------------------------
% Typography and spacing
% ------------------------------------------------------------
\setbeamerfont{title}{size=\LARGE,series=\bfseries}
\setbeamerfont{subtitle}{size=\large,series=\mdseries}
\setbeamerfont{author}{size=\normalsize,series=\mdseries}
\setbeamerfont{frametitle}{size=\Large,series=\bfseries}
\setbeamerfont{framesubtitle}{size=\normalsize,series=\mdseries}
\setbeamerfont{section in head/foot}{size=\scriptsize,series=\bfseries}
\setbeamerfont{subsection in head/foot}{size=\scriptsize,series=\mdseries}
\setbeamerfont{footline}{size=\scriptsize,series=\mdseries}
\setbeamerfont{block title}{series=\bfseries}

\setbeamersize{text margin left=1.15em,text margin right=1.15em}
\setlength{\parskip}{0.35em}

% ------------------------------------------------------------
% Bullets and links
% ------------------------------------------------------------
\setbeamercolor{item}{fg=TwilightDeepTeal}
\setbeamercolor{subitem}{fg=TwilightPurple}
\setbeamercolor{subsubitem}{fg=TwilightMulberry}
\setbeamertemplate{itemize item}{\raise0.2ex\hbox{\scriptsize$\blacktriangleright$}}
\setbeamertemplate{itemize subitem}{\raise0.15ex\hbox{\tiny$\blacktriangleright$}}

\setbeamercolor{progress active}{fg=TwilightMulberry}
\setbeamercolor{progress inactive}{fg=TwilightIvory}
\setbeamercolor{progress complete}{fg=TwilightDarkTeal}


% ------------------------------------------------------------
% Document Specific Data
% ------------------------------------------------------------
\title[ML FEM]{Multilevel Methods on Finite Elements}
\author[MB \& NK]{Melayne Barker and Neal Kuperman}
\institute{Portland State University}
\date{\today}
\definecolor{myactive}{HTML}{7A4A63} 
\definecolor{mycomplete}{HTML}{0D7C7B}
\definecolor{myinactive}{HTML}{F6F1EA}
\newcounter{progresssteps}
\setcounter{progresssteps}{18} % number of hats

\newcommand{\progressbar}{
\begin{tikzpicture}[remember picture,overlay]
    \foreach \i in {1,...,18} {
        % position along top
        \pgfmathsetmacro{\x}{(\i-1)*1.0}

        % determine if this hat should be active
        \pgfmathparse{\insertframenumber/\inserttotalframenumber*18 >= \i}
        \ifnum\pgfmathresult<1
            \colorlet{hatcolor}{myactive}
        \else
          \ifnum\pgfmathresult>1
            \colorlet{hatcolor}{myinactive}
          \else
            \colorlet{hatcolor}{mycomplete}
          \fi
        \fi

        % hat function
        \draw[fill=hatcolor,draw=hatcolor,fill opacity=0.25]
            ($(current page.north west)+(0.8+\x,-0.3)$)
            -- ++(0.35,0.25)
            -- ++(0.35,-0.25)
            -- cycle;
    }
\end{tikzpicture}
}
\begin{document}
\begin{centering}
\begin{frame}
  \titlepage
  \begin{tikzpicture}[remember picture, overlay]
    \node[
      anchor=south,
      inner sep=0pt
    ] at (current page.south) {
      \includegraphics[
        width=\paperwidth,
        height=0.64\paperheight
      ]{fem_title_graphic.png}
    };
  \end{tikzpicture}
\end{frame}

\section{The Problem}
\begin{frame}{Modeling Physics}
  \begin{columns}[T] 
\column{0.38\textwidth}
    \begin{itemize}
      \item Heat dissipating through a material.
      \item Problems include modeling how water flows from a source through porous materials. 
  \end{itemize}
\column{0.58\textwidth}
    \includegraphics[
        width=0.90\linewidth
        ]{heat-sink-model.png}
\end{columns}
\end{frame}

\begin{frame}{Modeling Physics}
  \begin{columns}[T] 
\column{0.38\textwidth}
    \begin{itemize}
    \item Problems include modeling how water flows from a source through porous materials.
    \item Heat dissipating through a material.
  \end{itemize}
\column{0.58\textwidth}
    \includegraphics[
          width=0.90\linewidth
    ]{discretization-heat-sink-model.png}
\end{columns}
\end{frame}

\begin{frame}{Modeling Physics}
  \begin{columns}[T] 
\column{0.38\textwidth}
    \begin{itemize}
    \item Problems include modeling how water flows from a source through porous materials.
    \item Heat dissipating through a material.
  \end{itemize}
\column{0.58\textwidth}
    \includegraphics[
          width=0.90\linewidth
    ]{temperature-field-heat-sink-model.png}
\end{columns}
\end{frame}

\section{The Approach}
\begin{frame}{PDEs and Finite Elements}
  To approximate the partial differential equation solutions, we use a basis of functions that behave in a manner we understand and that will simplify calculations.
  \begin{minipage}[t]{0.48\textwidth}
\centering
\includegraphics[width=\linewidth]{plot-using-discretization.png}
\\[0.5em]
\end{minipage}
\hfill
\begin{minipage}[t]{0.48\textwidth}
\centering
\includegraphics[width=\linewidth]{base-functions-overlap.png}
\\[0.5em]
\end{minipage}
\end{frame}

\begin{frame}{Key Elements of Finite Elements}
  \begin{enumerate}
    \item Start with a PDE
    \item Weak form
    \item Divergence theorem/integrate by parts
    \item Discretize the domain
    \item Represent the weak form as a discretization on the parts of the mesh
    \item The operators become matrices
    \item Solve with linear algebra
    \item Use norms to see how good your approximation is to toy problems where you know the solution
  \end{enumerate}
\end{frame}

\section{Our Goals}
\subsection{Adding Multilevel Methods}
\begin{frame}{Really Big Problems}
  \begin{columns}[T] 
  \column{0.78\textwidth}
    \includegraphics[
      width=0.90\linewidth
    ]{fluidspread.png}
  \column{0.28\textwidth}
    \begin{block}{Too big to run}
    Like Hillary said yesterday, some of these problems are absurdly large. We have to find ways to solve them when $x = A^(-1) b$
    \end{block} 
    
  \end{columns}
\end{frame}

\section{Our Methods}
\begin{frame}{NGSolve}
  \includegraphics[
      height=0.150\textheight
    ]{logo1.png}
  \includegraphics[
      height=0.750\textheight
    ]{ngsolve.png}
\end{frame}

\subsection{Our Constructed Framework}
\begin{frame}{A Hierarchy of Finite Element Spaces}
   \includegraphics[
      height=0.750\textheight
    ]{multigrid.png}
\end{frame}

\begin{frame}{Setting up the spaces}
  At each level, we have the data storage, solving tools, and projection matrices available.
  \includegraphics[
    height=0.75\textheight
  ]{level-parts.png}
\end{frame}

\section{Our Results}
\begin{frame}{Gauss Siedel Smoothing}
  \begin{columns}[T] 
\column{0.38\textwidth}
    \begin{example}{Two level method}\\
      Use Gauss Siedel to smooth the matrix, then
      \begin{align}
        x_{k+1} &= b - M^{-1}Ax_k
      \end{align}
      fine ndof=1961, coarse ndof=511,  $P$ shape 1961x511\\
      $\left\| e \right\|_E $ (after section 1 pre-smooth): 1.790578e+00\\
      $\left\| e \right\|_E$ after coarse correction:      1.126177e-01\\
      $\left\| e \right\|$ L2 error:                      9.820784e-04
    \end{example}

\column{0.58\textwidth}
    \includegraphics[
          width=0.90\linewidth
        ]{smooth1.png}

\end{columns}
\end{frame}

\begin{frame}{Gauss Siedel Smoothing}
  \begin{columns}[T] 
\column{0.38\textwidth}
    \begin{example}{Single-level smoother}\\
      fine ndof=1961, coarse ndof=511,  $P$ shape 1961x511\\
      $\left\| e \right\|_E $ (after section 1 pre-smooth): 1.790578e+00\\
      $\left\| e \right\|_E$ after coarse correction:      1.126177e-01\\
      $\left\| e \right\|$ L2 error:                      9.820784e-04
    \end{example}

\column{0.58\textwidth}
    \includegraphics[
          width=0.90\linewidth
        ]{smooth2.png}

\end{columns}
\end{frame}

\begin{frame}{Gauss Siedel Smoothing}
  \begin{columns}[T] 
\column{0.38\textwidth}
    \begin{example}{Single-level smoother}\\
      fine ndof=1961, coarse ndof=511,  $P$ shape 1961x511\\
      $\left\| e \right\|_E $ (after section 1 pre-smooth): 1.790578e+00\\
      $\left\| e \right\|_E$ after coarse correction:      1.126177e-01\\
      $\left\| e \right\|$ L2 error:                      9.820784e-04
    \end{example}

\column{0.58\textwidth}
    \includegraphics[
          width=0.90\linewidth
        ]{output1.png}

\end{columns}
\end{frame}

\begin{frame}{Solving V-cycle with smoothing and coarse solve}
  \includegraphics[
          width=\textwidth
        ]{vcycle.png}
\end{frame}

\begin{frame}{Convergence of v-cycle}
  \begin{columns}[T] 
\column{0.48\textwidth}
   \includegraphics[
          width=0.90\linewidth
        ]{conv_plot2.png}
        \column{0.52\textwidth}
        \includegraphics[
         height=0.85\textheight
        ]{conv_plot3.png}
  \end{columns}
\end{frame}

\begin{frame}{Now we use it as a preconditioner}
  \begin{alert}{Success!}
  We made it so that we run PCG and when it asks for a preconditioner, it gets our two-level method
  \end{alert}
  PCG (NGSolve CG + your V-cycle): 1684 iterations
\end{frame}

\section{Conclusion}
\begin{frame}{Conclusion}
  \begin{itemize}
    \item We demonstrated the smoothing effects of Gauss Siedel on highly oscillatory components.
    \item We implemented a two-level method with GS smoothing and a coarse direct solve.
    \item We built a framework from which future multilevel test on NGSolve finite element spaces may be performed and tested it to demonstrate convergence.
    \item We implemented the method above as the preconditioner for a preconditioned conjugate gradient solver.
  \end{itemize}
\end{frame}
\end{centering}

\section{References}
\begin{frame}
  \begin{itemize}
    \item M. la Cour Christensen, U. Villa, A. P. Engsig-Karup, and P. S. Vassilevski, "Numerical Multilevel Upscaling for Incompressible Flow in Reservoir Simulation: An Element-Based Algebraic Multigrid (AMGe) Approach," SIAM J. Sci. Comput., vol. 39, no. 1, pp. B102–B137, 2017.
    \item P. S. Vassilevski, "Multilevel Block Factorization Preconditioners, Matrix-based Analysis and Algorithms for Solving Finite Element Equations", Springer, New York, 2008. 514 p.
    \item J. S. Ovall, Numerical Mathematics. Philadelphia, PA: SIAM, 2024. DOI: 10.1137/1.9781611978070.
    \item J. N. Reddy, N. K. Anand, and P. Roy, Finite Element and Finite Volume Methods for Heat Transfer and Fluid Dynamics. Cambridge University Press, Cambridge, UK, 2023. DOI: 10.1017/9781009275453.
    \item Schöberl, J., et al. NGSolve Interactive Tutorials (i-Tutorials). Netgen/NGSolve Documentation. Available at: https://docu.ngsolve.org/latest/i-tutorials/ (accessed June 2026).
    \item Some finite element images came from https://www.comsol.com/multiphysics/fea-software
    \item Dilip Krishnan, Richard Szeliski
"Multigrid and Multilevel Preconditioners for Computational
Photography" https://www.microsoft.com/en-us/research/wp-content/uploads/2011/12/Krishanan-SGA11-Supp.pdf  
\end{itemize}
\end{frame}
\end{document}
